\documentclass[11pt,a4paper]{article}

% ---------------------------------------------------------
% PREAMBLE (stable, warning-free, Overleaf-safe)
% ---------------------------------------------------------
\usepackage[utf8]{inputenc}
\usepackage[T1]{fontenc}
\usepackage{lmodern}
\usepackage{amsmath, amssymb, amsfonts}
\usepackage{bm}
\usepackage{geometry}
\usepackage{graphicx}
\usepackage{hyperref}
\usepackage{cite}
\usepackage{authblk}
\usepackage{physics}
\usepackage{mathtools}
\usepackage{enumitem}
\usepackage{tensor}

\geometry{margin=2.4cm}

% ---------------------------------------------------------
% TITLE
% ---------------------------------------------------------
\title{\textbf{RDCC Summary Paper V2.0:\\
The RDCC Inference Framework: Geometry, Topology, Likelihoods,\\
and Euclid-Level Forecasts}}

\author[1]{Michael Lehmann}
\affil[1]{\small Independent Researcher, Relaxation-Driven Cyclic Cosmology (RDCC) Framework\\
\texttt{mi.lehmann@gmx.de}}

\date{\small June 2026}

\begin{document}
\maketitle

% ========================================================
% ABSTRACT
% ========================================================

\begin{abstract}
The Relaxation-Driven Cyclic Cosmology (RDCC) modifies the gravitational potential
through a scale- and redshift-dependent relaxation kernel $\mathcal{R}(k,z;\alpha)$.
These modifications induce geometric, morphological, and topological signatures in
weak-lensing observables that require an inference framework capable of capturing
structure beyond two-point statistics. Over the RDCC Companion Series (LXXX--C),
we developed a complete inference ecosystem combining geometric response theory,
persistent homology, transport-based topology, Minkowski functionals, emulator
pipelines, multi-probe likelihoods, Fisher geometry, and degeneracy analysis.

This Summary Paper synthesizes these developments into a unified, survey-ready
framework. We present: (i) the geometric and topological structure of RDCC
signatures, (ii) the construction of transport-based and morphological summary
statistics, (iii) the architecture of the RDCC likelihood, (iv) the Fisher geometry
and information-flow structure underlying RDCC inference, (v) optimal RDCC probes
and survey strategies, and (vi) Euclid-level forecasts for the relaxation parameter
$\alpha$. This document serves as the central reference for the RDCC inference
ecosystem and provides a coherent foundation for future observational applications.
\end{abstract}

% ========================================================
% 1 INTRODUCTION
% ========================================================

\section{Introduction}

The Relaxation-Driven Cyclic Cosmology (RDCC) introduces a dynamical relaxation
mechanism that modifies the gravitational potential through a single parameter
$\alpha$. These modifications produce scale-dependent distortions in the lensing
potential that propagate into geometric, morphological, and topological features
of the weak-lensing convergence field. Capturing these signatures requires an
inference framework that extends beyond two-point statistics and incorporates
non-Gaussian structure, geometric sensitivity, and topological information.

Over the RDCC Companion Series (LXXX--C), we developed such a framework. It
combines geometric response fields, persistent homology, transport-based topology,
Minkowski functionals, emulator pipelines, multi-probe likelihoods, Fisher
geometry, and degeneracy analysis into a coherent and modular ecosystem.

The RDCC inference framework rests on four conceptual pillars:

\begin{enumerate}[leftmargin=1.2cm]
    \item \textbf{Geometric Signatures:}  
    RDCC induces scale-dependent distortions in the lensing potential, encoded in
    geometric response fields, Fisher channels, and curvature-aware parameter-space
    metrics.

    \item \textbf{Topological Structure:}  
    RDCC modifies the birth--death structure of persistent homology, the morphology
    of peaks and voids, and the transport geometry of topological summaries.

    \item \textbf{Likelihood Architecture:}  
    A modular likelihood integrates transport-based statistics, Minkowski
    functionals, peak/void statistics, power spectra, and neural-compressed
    summaries, supported by emulator pipelines and simulation-based calibration.

    \item \textbf{Forecasts and Model Selection:}  
    Euclid-level forecasts, Fisher-channel analysis, and evidence decomposition
    quantify the detectability of RDCC and its distinguishability from
    $\Lambda$CDM.
\end{enumerate}

This Summary Paper synthesizes these components into a unified narrative. Rather
than reproducing the detailed derivations contained in the individual Companions,
we highlight the structural logic of the RDCC inference framework, the interplay
between geometry and topology, and the methodological innovations that enable
robust and high-fidelity parameter estimation.

% ========================================================
% 2 RDCC MODEL AND GEOMETRIC STRUCTURE
% ========================================================

\section{The RDCC Model and Its Geometric Structure}

RDCC modifies the gravitational potential through a relaxation kernel
\begin{equation}
    \Phi_{\mathrm{RDCC}}(k,z)
    =
    \Phi_{\Lambda\mathrm{CDM}}(k,z)\,
    \mathcal{R}(k,z;\alpha),
\end{equation}
where $\mathcal{R}(k,z;\alpha)$ introduces scale-dependent deviations that become
relevant at $k \gtrsim 0.1\,h\,\mathrm{Mpc}^{-1}$. These deviations propagate into
weak-lensing observables and generate geometric and topological signatures that
form the basis of RDCC inference.

\subsection{Geometric Response Fields}

For an observable $O$, the geometric response to $\alpha$ is defined as
\begin{equation}
    \mathcal{R}_{O}
    =
    \pdv{O}{\alpha},
\end{equation}
which quantifies the local deformation induced by RDCC. These response fields
capture peak sharpening, void broadening, filament-void boundary shifts, and
curvature changes in the lensing potential.

\subsection{Fisher Geometry}

The RDCC parameter space inherits a Riemannian structure from the Fisher
information metric
\begin{equation}
    g_{\alpha\alpha}
    =
    \mathbb{E}\!\left[
        \left(
            \pdv{\log \mathcal{L}}{\alpha}
        \right)^{2}
    \right],
\end{equation}
which encodes the curvature of the likelihood surface. Fisher geometry reveals
RDCC-sensitive directions, degeneracy manifolds, and natural coordinates for
efficient inference.

% ========================================================
% COMPANION ROADMAP
% ========================================================

\begin{table}[t]
\centering
\renewcommand{\arraystretch}{1.25}
\begin{tabular}{p{2.2cm} p{6.2cm} p{6.2cm}}
\hline
\textbf{Companion} & \textbf{Content} & \textbf{Role in RDCC Framework} \\
\hline

LXXX--LXXXII &
Geometric RDCC signatures, response fields, scale-dependent distortions &
Defines geometric foundations; identifies RDCC-sensitive scales and response patterns \\

LXXXIII--LXXXVI &
Topological statistics, persistent homology, transport geometry &
Establishes non-Gaussian RDCC signatures; provides robust topology-based probes \\

LXXXVII--LXXXVIII &
Minkowski functionals, morphological descriptors &
Captures RDCC-induced shape and connectivity changes in convergence maps \\

LXXXIX--XC &
Emulator pipelines, simulation grids, interpolation architecture &
Enables fast likelihood evaluation; propagates simulation uncertainty \\

XCI--XCII &
Forecasts, Euclid-like constraints, multi-probe combinations &
Quantifies detectability of $\alpha$; identifies optimal probe combinations \\

XCIII--XCIV &
Fisher geometry, curvature, natural coordinates &
Defines RDCC parameter-space geometry; identifies sensitivity directions \\

XCV--XCVI &
Information tensor, Fisher channels, information flow &
Decomposes RDCC information; identifies dominant channels and null directions \\

XCVII--XCVIII &
Degeneracy structure, null space, RDCC--$\Lambda$CDM alignment &
Maps degeneracy manifolds; identifies blind modes and orthogonal probes \\

XCIX &
Optimal probes, survey strategies, efficiency function &
Ranks probes; provides survey recommendations for RDCC detection \\

C &
Benchmark Suite, validation, stress tests, reproducibility framework &
Ensures robustness, emulator validation, and cross-probe consistency \\
\hline
\end{tabular}
\caption{Roadmap of the RDCC Companion Series. Each Companion contributes a 
specific conceptual or technical component to the RDCC inference ecosystem.}
\end{table}

% ========================================================
% 3 TOPOLOGICAL AND TRANSPORT-BASED RDCC SIGNATURES
% ========================================================

\section{Topological and Transport-Based RDCC Signatures}

Weak-lensing observables contain rich geometric and topological information that
extends beyond two-point statistics. RDCC induces characteristic distortions in
the morphology of the convergence field, the connectivity of the cosmic web, and
the birth--death structure of persistent homology. These effects motivate the use
of topological and transport-based summary statistics, which provide
noise-resistant and non-Gaussian probes of RDCC signatures.

% --------------------------------------------------------
\subsection{Persistent Homology and Birth--Death Structure}

Persistent homology characterizes the topology of the convergence field by
tracking the evolution of connected components, loops, and voids across
thresholds. For a filtration parameter $\nu$, the persistence diagram
$\mathcal{D}(\nu)$ encodes the birth and death of topological features.

RDCC modifies the persistence structure through:

\begin{itemize}
    \item enhanced persistence of filamentary loops,
    \item earlier birth of small-scale peaks,
    \item delayed death of large voids,
    \item scale-dependent shifts in birth--death trajectories.
\end{itemize}

These effects are captured by the RDCC topological gradient
\begin{equation}
    \mathcal{G}_{\mathrm{topo}}
    =
    \pdv{\mathcal{T}}{\alpha},
\end{equation}
where $\mathcal{T}$ denotes a transport-based topological statistic.

% --------------------------------------------------------
\subsection{Transport Geometry of Topological Statistics}

Transport-based statistics quantify the geometric displacement between
persistence diagrams. Let $\mathcal{D}_{\mathrm{RDCC}}$ and
$\mathcal{D}_{\Lambda\mathrm{CDM}}$ denote diagrams under RDCC and
$\Lambda$CDM. Their quadratic optimal-transport distance is
\begin{equation}
    W_{2}\!\left(
        \mathcal{D}_{\mathrm{RDCC}},
        \mathcal{D}_{\Lambda\mathrm{CDM}}
    \right),
\end{equation}
which measures the minimal geometric deformation required to transform one
diagram into the other.

RDCC induces coherent transport flows that reflect:

\begin{itemize}
    \item peak migration toward higher persistence,
    \item void expansion and boundary smoothing,
    \item filament thickening or thinning,
    \item redistribution of topological mass across scales.
\end{itemize}

These flows provide a geometric representation of RDCC-induced topological
deformations.

% --------------------------------------------------------
\subsection{Minkowski Functionals and Morphological Signatures}

Minkowski functionals $V_{i}(\nu)$ capture the morphology of the convergence
field across thresholds. RDCC modifies:

\begin{itemize}
    \item the area functional $V_{0}$ through small-scale suppression,
    \item the boundary length $V_{1}$ via filament--void transitions,
    \item the Euler characteristic $V_{2}$ through changes in peak/void counts.
\end{itemize}

These morphological signatures complement persistent homology and provide
summary statistics that are computationally efficient and robust to noise.

% --------------------------------------------------------
\subsection{Joint Topological--Geometric Structure}

Topological and geometric signatures are not independent. RDCC induces coherent
patterns across:

\begin{itemize}
    \item geometric response fields,
    \item transport flows in persistence space,
    \item Minkowski functional derivatives,
    \item Fisher-channel alignments.
\end{itemize}

This joint structure is essential for constructing multi-probe likelihoods and
for identifying optimal RDCC probes.

% --------------------------------------------------------
\subsection{Summary}

Topological and transport-based statistics provide a powerful and noise-resistant
characterization of RDCC signatures. They capture the non-Gaussian structure of
the convergence field, reveal scale-dependent deformations induced by the
relaxation mechanism, and form a central component of the RDCC inference
framework. Their integration with geometric sensitivity maps and Fisher geometry
enables a comprehensive and robust analysis of RDCC effects.

% ========================================================
% 4 THE RDCC LIKELIHOOD ARCHITECTURE
% ========================================================

\section{The RDCC Likelihood Architecture}

The RDCC inference framework requires a likelihood capable of integrating
non-Gaussian summary statistics, topological descriptors, geometric sensitivity
patterns, and emulator-based predictions. The resulting architecture is modular,
hierarchical, and designed to combine heterogeneous probes while maintaining
statistical consistency and computational efficiency.

% --------------------------------------------------------
\subsection{Multi-Probe Structure}

Let $\mathcal{S}$ denote the full set of RDCC summary statistics, including:

\begin{itemize}
    \item power spectra and correlation functions,
    \item Minkowski functionals,
    \item transport-based topological statistics,
    \item peak and void counts,
    \item neural-compressed summaries.
\end{itemize}

The joint likelihood factorizes as
\begin{equation}
    \mathcal{L}(\mathcal{S}\,|\,\alpha,\theta)
    =
    \prod_{i}
    \mathcal{L}_{i}(S_{i}\,|\,\alpha,\theta),
\end{equation}
where $\theta$ denotes the $\Lambda$CDM parameters. Each component likelihood
$\mathcal{L}_{i}$ is calibrated using simulation-based inference and validated
through cross-probe consistency tests.

% --------------------------------------------------------
\subsection{Emulator-Based Predictions}

High-fidelity RDCC simulations are computationally expensive. To enable efficient
likelihood evaluation, we employ emulator pipelines trained on simulation grids
and validated using the RDCC Benchmark Suite (Companion C).

For each summary statistic $S_{i}$, the emulator predicts
\begin{equation}
    S_{i}(\alpha,\theta),
\end{equation}
with uncertainty propagated into the likelihood via
\begin{equation}
    \Sigma_{i}
    =
    \Sigma_{\mathrm{sim}}
    +
    \Sigma_{\mathrm{emul}}.
\end{equation}

This ensures that emulator uncertainty does not bias inference or artificially
inflate evidence contrast.

% --------------------------------------------------------
\subsection{Transport-Based Likelihood Components}

Transport-based statistics require likelihoods that respect the geometry of
persistence diagrams. For a quadratic optimal-transport distance $W_{2}$, we
define
\begin{equation}
    \mathcal{L}_{\mathrm{topo}}
    \propto
    \exp\!\left[
        -\frac{
            W_{2}\!\left(
                \mathcal{D}_{\mathrm{data}},
                \mathcal{D}(\alpha)
            \right)
        }{
            2\,\sigma_{W_{2}}^{2}
        }
    \right],
\end{equation}
which captures geometric deformations in topological space and is robust to noise
and sampling variance.

% --------------------------------------------------------
\subsection{Neural Compression and Likelihood-Free Components}

For high-dimensional statistics, we employ neural compression to construct
compressed summaries $z$ that retain maximal information about $\alpha$:
\begin{equation}
    z = f_{\mathrm{enc}}(\mathcal{S}),
\end{equation}
where $f_{\mathrm{enc}}$ is trained to maximize Fisher information or minimize
likelihood-free loss functions.

The compressed likelihood is then
\begin{equation}
    \mathcal{L}(z\,|\,\alpha)
    =
    \mathcal{N}\!\left(
        z;\,
        \mu(\alpha),\,
        \Sigma_{z}
    \right),
\end{equation}
with $\Sigma_{z}$ estimated via simulation-based calibration.

% --------------------------------------------------------
\subsection{Cross-Probe Covariance Structure}

The full covariance matrix is
\begin{equation}
    C
    =
    C_{\mathrm{stat}}
    +
    C_{\mathrm{sys}}
    +
    C_{\mathrm{cross}},
\end{equation}
where $C_{\mathrm{cross}}$ captures correlations between geometric, topological,
and morphological probes.

RDCC induces characteristic cross-covariance patterns, including:

\begin{itemize}
    \item alignment between geometric response fields and topological flows,
    \item covariance suppression in degeneracy-aligned directions,
    \item enhanced cross-probe coherence at RDCC-sensitive scales.
\end{itemize}

% --------------------------------------------------------
\subsection{Summary}

The RDCC likelihood architecture integrates geometric, topological, and
compressed statistics into a unified multi-probe framework. Emulator-based
predictions, transport-aware likelihoods, and cross-probe covariance modeling
ensure robustness and flexibility. This architecture forms the computational core
of RDCC inference and underpins the forecasts and model-selection results
presented in later sections.

% ========================================================
% 5 FISHER GEOMETRY, INFORMATION FLOW, AND DEGENERACY STRUCTURE
% ========================================================

\section{Fisher Geometry, Information Flow, and Degeneracy Structure}

The RDCC inference framework is governed by a geometric structure induced by the
Fisher information metric, which encodes the curvature of the likelihood surface
and determines the efficiency of parameter estimation. Combined with observable
gradients and topological response fields, the Fisher geometry defines the
information flow from weak-lensing observables to RDCC parameter constraints.

% --------------------------------------------------------
\subsection{Fisher Information and the RDCC Metric}

For a parameter vector $(\alpha,\theta)$, the Fisher information matrix is
\begin{equation}
    G_{ij}
    =
    \mathbb{E}\!\left[
        \pdv{\log \mathcal{L}}{\theta_{i}}
        \pdv{\log \mathcal{L}}{\theta_{j}}
    \right],
\end{equation}
where $\theta$ denotes the $\Lambda$CDM parameters. The Fisher metric defines a
Riemannian geometry on parameter space, with geodesics corresponding to
directions of minimal information loss.

For RDCC, the Fisher metric reveals:

\begin{itemize}
    \item strong curvature in RDCC-sensitive directions,
    \item flat regions corresponding to degeneracy manifolds,
    \item natural coordinates that diagonalize information channels,
    \item geometric pathways linking RDCC and $\Lambda$CDM variations.
\end{itemize}

These geometric features guide the construction of efficient samplers and inform
the design of optimal probes.

% --------------------------------------------------------
\subsection{Observable Gradients and Information Flow}

Let $S_{i}$ denote a summary statistic. Its RDCC gradient is
\begin{equation}
    \mathcal{G}_{i}
    =
    \pdv{S_{i}}{\alpha}.
\end{equation}

The information tensor introduced in Companion XCVII is
\begin{equation}
    I_{ij}
    =
    \mathbb{E}\!\left[
        \mathcal{G}_{i}\,
        \mathcal{G}_{j}
    \right],
\end{equation}
which generalizes the Fisher metric to the space of summary statistics.

Diagonalizing $I$ yields Fisher channels:
\begin{equation}
    I\,v_{k}
    =
    \lambda_{k}\,v_{k}.
\end{equation}

Each channel $v_{k}$ represents:

\begin{itemize}
    \item a direction of coherent RDCC response,
    \item a geometric mode of information flow,
    \item a contribution $\lambda_{k}$ to the total Fisher information,
    \item a pathway linking observables to evidence.
\end{itemize}

Large eigenvalues correspond to dominant RDCC information pathways, while small
eigenvalues indicate suppressed or degenerate directions.

% --------------------------------------------------------
\subsection{Degeneracy Structure and Null Directions}

Not all directions in summary-statistic space carry RDCC information. The RDCC
Null Space is defined as
\begin{equation}
    \mathcal{N}
    =
    \left\{
        v \,\middle|\,
        v^{T} I v = 0
    \right\}.
\end{equation}

Vectors in $\mathcal{N}$ correspond to:

\begin{itemize}
    \item blind modes with zero RDCC sensitivity,
    \item combinations of observables aligned with $\Lambda$CDM variations,
    \item directions dominated by noise or cosmic variance,
    \item degeneracy manifolds in Fisher geometry.
\end{itemize}

These null directions play a central role in probe optimization and survey
design, as they identify observables that do not contribute to RDCC inference.

% --------------------------------------------------------
\subsection{Geometric Interpretation of Degeneracy Manifolds}

In Fisher geometry, degeneracy manifolds satisfy
\begin{equation}
    g_{\alpha\alpha}\,d\alpha^{2}
    =
    g_{\theta\theta}\,d\theta^{2},
\end{equation}
indicating that variations in $\alpha$ can be mimicked by variations in
$\Lambda$CDM parameters.

These manifolds represent:

\begin{itemize}
    \item curved surfaces of equal RDCC sensitivity,
    \item geodesic alignments between RDCC and $\Lambda$CDM,
    \item regions of suppressed evidence contrast,
    \item pathways of minimal distinguishability.
\end{itemize}

Understanding these manifolds is essential for breaking degeneracies and
designing optimal RDCC probes.

% --------------------------------------------------------
\subsection{Summary}

Fisher geometry, information flow, and degeneracy structure form the
mathematical backbone of RDCC inference. Observable gradients propagate through
the information tensor to define Fisher channels, while degeneracy manifolds
identify directions of suppressed sensitivity. Together, these structures guide
the construction of efficient likelihoods, optimal probes, and robust forecasting
pipelines.

% ========================================================
% 6 OPTIMAL PROBES AND SURVEY STRATEGIES
% ========================================================

\section{Optimal Probes and Survey Strategies}

The RDCC inference framework integrates geometric sensitivity, topological
structure, and Fisher-channel analysis to identify the most informative
observables for constraining the relaxation parameter $\alpha$. This section
summarizes the optimal RDCC probes and the survey strategies that maximize
detectability and minimize degeneracy with $\Lambda$CDM.

% --------------------------------------------------------
\subsection{The RDCC Probe Efficiency Function}

To quantify the information content of a summary statistic $S$, we define the
RDCC Probe Efficiency Function
\begin{equation}
    \varepsilon(S)
    =
    \left(
        1 - D(S)
    \right)
    \frac{
        \mathcal{I}(S)
    }{
        \mathrm{Var}(S)
    },
\end{equation}
where $D(S)$ is the degeneracy factor derived from the RDCC Null Space
(Companion XCVIII), and $\mathcal{I}(S)$ denotes the Fisher information
contributed by $S$.

High values of $\varepsilon(S)$ indicate probes that:

\begin{itemize}
    \item respond strongly to RDCC-induced distortions,
    \item break degeneracies with $\Lambda$CDM parameters,
    \item align with dominant Fisher channels,
    \item exhibit stable topological or geometric signatures.
\end{itemize}

This efficiency function provides a unified ranking of RDCC probes.

% --------------------------------------------------------
\subsection{Ranking RDCC Probes}

Evaluating $\varepsilon(S)$ across geometric, topological, and morphological
statistics yields a consistent hierarchy:

\begin{itemize}
    \item \textbf{Transport-based topological statistics}  
    provide the strongest RDCC sensitivity, particularly at intermediate
    persistence scales.

    \item \textbf{Minkowski functionals}  
    capture morphological distortions induced by RDCC and complement topological
    probes.

    \item \textbf{Peak and void statistics}  
    are highly sensitive to relaxation-driven shifts in the potential.

    \item \textbf{Power spectra}  
    contribute primarily at large scales but are partially degenerate with
    $\Lambda$CDM variations.

    \item \textbf{Neural-compressed summaries}  
    efficiently combine multiple probes and retain high Fisher information.
\end{itemize}

This ranking reflects the multi-scale and non-Gaussian nature of RDCC
signatures.

% --------------------------------------------------------
\subsection{Fisher-Channel Alignment}

The geometric alignment of a probe $S$ with the dominant Fisher channel
$v_{\max}$ is
\begin{equation}
    A(S)
    =
    \frac{
        \left\langle
            \pdv{S}{\alpha},
            v_{\max}
        \right\rangle
    }{
        \norm{\pdv{S}{\alpha}}\,\norm{v_{\max}}
    }.
\end{equation}

Probes with high alignment:

\begin{itemize}
    \item efficiently trace the dominant RDCC information pathway,
    \item minimize curvature-induced sampling inefficiencies,
    \item reduce the impact of degeneracy manifolds,
    \item provide stable likelihood gradients.
\end{itemize}

Transport-based topological statistics exhibit the highest alignment, followed
by Minkowski functionals and peak statistics.

% --------------------------------------------------------
\subsection{Breaking Degeneracies with \texorpdfstring{$\Lambda$CDM}{ΛCDM}}

Degeneracy manifolds identified in Companion XCVIII reveal directions in
summary-statistic space where RDCC variations mimic $\Lambda$CDM parameter
shifts. Optimal probes are those that:

\begin{itemize}
    \item have minimal projection onto the RDCC Null Space,
    \item exhibit strong orthogonality to $\Lambda$CDM-aligned subspaces,
    \item amplify RDCC-specific geometric or topological distortions.
\end{itemize}

Topological probes are particularly effective at breaking degeneracies due to
their sensitivity to non-Gaussian structure.

% --------------------------------------------------------
\subsection{Survey Strategies for Maximizing RDCC Detectability}

The RDCC probe hierarchy informs survey strategies that maximize sensitivity to
$\alpha$. Key recommendations include:

\begin{itemize}
    \item \textbf{Angular resolution:}  
    prioritize intermediate scales where RDCC-induced topological distortions are
    strongest.

    \item \textbf{Tomographic binning:}  
    allocate finer redshift resolution to bins with high Fisher-channel
    alignment.

    \item \textbf{Noise mitigation:}  
    emphasize probes with stable topological persistence under shape noise.

    \item \textbf{Cross-probe integration:}  
    combine geometric and topological probes to suppress degeneracy directions.

    \item \textbf{Emulator fidelity:}  
    ensure high accuracy for probes with large $\varepsilon(S)$ to avoid
    information loss.
\end{itemize}

These strategies are compatible with Euclid-like survey specifications and
provide a roadmap for future RDCC analyses.

% --------------------------------------------------------
\subsection{Summary}

Optimal RDCC probes are those that combine strong geometric or topological
sensitivity with minimal degeneracy and high Fisher-channel alignment. The
resulting probe hierarchy and survey strategies maximize RDCC detectability and
provide a principled foundation for high-precision cosmological inference.

% ========================================================
% 7 FORECASTS AND MODEL SELECTION
% ========================================================

\section{Forecasts and Model Selection}

The RDCC inference framework enables quantitative forecasts for the relaxation
parameter $\alpha$ and provides a principled basis for model comparison with
$\Lambda$CDM. These forecasts combine geometric sensitivity, topological
statistics, emulator-based predictions, and the multi-probe likelihood
architecture described in previous sections.

% --------------------------------------------------------
\subsection{Euclid-Level Constraints on \texorpdfstring{$\alpha$}{α}}

Using the full RDCC likelihood, including transport-based topology, Minkowski
functionals, peak statistics, and neural-compressed summaries, we obtain
forecasted constraints of the form
\begin{equation}
    \sigma(\alpha)
    \sim
    \mathcal{O}(10^{-2}),
\end{equation}
with the precise value depending on the probe combination and tomographic
binning.

The strongest constraints arise from:

\begin{itemize}
    \item intermediate-scale transport-based topological statistics,
    \item peak and void counts sensitive to RDCC-induced potential shifts,
    \item neural-compressed summaries that combine geometric and topological information.
\end{itemize}

Power spectra alone provide significantly weaker constraints due to partial
degeneracy with $\Lambda$CDM parameters.

% --------------------------------------------------------
\subsection{Fisher-Channel Contributions to Forecasts}

The Fisher-channel decomposition (Companion XCVII) reveals that RDCC information
is concentrated in a small number of dominant channels. The leading channel
typically accounts for
\begin{equation}
    \sim 70\%
\end{equation}
of the total Fisher information, with the remaining channels capturing
subdominant geometric and topological features.

This structure explains:

\begin{itemize}
    \item the efficiency of neural compression,
    \item the robustness of topological probes,
    \item the limited contribution of high-dimensional statistics,
    \item the sensitivity of forecasts to degeneracy directions.
\end{itemize}

% --------------------------------------------------------
\subsection{Degeneracy Structure and Forecast Stability}

Degeneracy manifolds identified in Companion XCVIII reveal directions in
parameter space where RDCC variations mimic $\Lambda$CDM shifts. These
degeneracies affect forecast stability by:

\begin{itemize}
    \item reducing sensitivity in null directions,
    \item suppressing evidence contrast in aligned subspaces,
    \item amplifying the importance of topological probes,
    \item increasing the role of cross-probe covariance.
\end{itemize}

Forecasts that rely solely on geometric probes are more susceptible to these
degeneracies, while topological probes provide strong orthogonality to
$\Lambda$CDM-aligned directions.

% --------------------------------------------------------
\subsection{Evidence Decomposition and Model Selection}

The RDCC evidence can be decomposed into Fisher-channel contributions:
\begin{equation}
    \Delta \log Z_{k}
    =
    \frac{1}{2}\,\lambda_{k},
\end{equation}
where $\lambda_{k}$ is the eigenvalue of the $k$-th Fisher channel. Summing over
channels yields the total evidence contrast between RDCC and $\Lambda$CDM.

Topological probes contribute disproportionately to $\Delta \log Z$, reflecting
their sensitivity to non-Gaussian RDCC signatures. Geometric probes contribute
primarily through the leading Fisher channel, while power spectra contribute
weakly due to degeneracy with $\Lambda$CDM.

% --------------------------------------------------------
\subsection{Combined Forecasts and Model-Selection Prospects}

Combining all probes yields:

\begin{itemize}
    \item tight constraints on $\alpha$,
    \item strong evidence contrast for moderate RDCC amplitudes,
    \item robustness to noise and systematics,
    \item stability under cross-probe consistency tests.
\end{itemize}

The RDCC inference framework therefore provides a powerful and flexible approach
to testing relaxation-driven cosmology with upcoming weak-lensing surveys.

% --------------------------------------------------------
\subsection{Summary}

Forecasts and model-selection analyses demonstrate that RDCC induces detectable
geometric and topological signatures in weak-lensing observables. Topological
probes, Fisher-channel structure, and degeneracy analysis play central roles in
achieving high sensitivity and robust model discrimination. These results
establish RDCC as a testable extension of $\Lambda$CDM and motivate future
observational applications.

% ========================================================
% 8 CONCLUSION
% ========================================================

\section{Conclusion}

The Relaxation-Driven Cyclic Cosmology (RDCC) introduces a scale-dependent
relaxation mechanism that produces geometric and topological signatures in
weak-lensing observables. Over the course of the RDCC Companion Series
(LXXX--C), we developed a complete inference framework capable of capturing
these signatures across multiple statistical domains. This Summary Paper
synthesized the conceptual foundations, methodological components, and
forecasting results of that framework.

RDCC induces coherent deformations in the gravitational potential, which
propagate into geometric response fields, Fisher-channel structure, and
persistent topological features. These signatures motivate the use of
transport-based topology, Minkowski functionals, peak and void statistics, and
neural-compressed summaries. The resulting multi-probe likelihood architecture
integrates these heterogeneous observables into a unified inference pipeline,
supported by emulator-based predictions and validated through the RDCC
Benchmark Suite.

Fisher geometry and information-flow analysis reveal that RDCC information is
concentrated in a small number of dominant channels, while degeneracy manifolds
identify directions of suppressed sensitivity and partial alignment with
$\Lambda$CDM variations. Optimal probes are those that combine strong geometric
or topological response with minimal degeneracy and high Fisher-channel
alignment. These probes form the basis of Euclid-level forecasts, which
demonstrate that RDCC induces detectable signatures and that the relaxation
parameter $\alpha$ can be constrained with high precision.

The RDCC inference framework is modular, extensible, and designed for
application to upcoming weak-lensing surveys. Its combination of geometric
analysis, topological structure, and simulation-based inference provides a robust
foundation for testing relaxation-driven cosmology. Future work will extend this
framework to cross-correlations with galaxy clustering, multi-wavelength probes,
and joint analyses with dynamical tracers. The RDCC Companion Series serves as a
technical backbone for these developments and ensures reproducibility,
consistency, and methodological clarity across the entire RDCC ecosystem.

\end{document}